\providecommand{\vf}[1]{\textcolor{blue}{#1}}
\chapter{Literature Review}\label{ch:lit}

This chapter reviews attacks that reconstruct an input from a representation, and it orders them by what the adversary must hold.
Section~\ref{sec:lit:split} covers the hidden states of a split model, Section~\ref{sec:lit:pooled} covers pooled text embeddings, Section~\ref{sec:lit:align} covers the alignment of one model's representations with another's, and Section~\ref{sec:lit:defense} covers defenses.
Section~\ref{sec:lit:position} places the dissertation among them.

\section{Inversion of Hidden States in Split Inference}\label{sec:lit:split}

Early work on the inversion of intermediate representations studied images.
He et al.~\cite{he2019model} split a convolutional network between a client and a server and recover the input image from the activations at the split layer, in three settings: with the model's weights, with queries to the model, and with neither.
Pasquini et al.~\cite{pasquini2021unleashing} attack split learning, the training counterpart: a malicious server steers the client's network during training until its activations can be inverted, with a public dataset and no knowledge of the client's data.
The attack is active and acts during training, whereas the adversary of this dissertation is passive and observes inference only.

\paragraph{Attacks with the model's weights.} Several attacks on language models assume a white-box adversary, one that holds the weights of the blocks that produced the hidden states.
Qu et al.~\cite{qu2025prompt} optimize an input embedding until its hidden states match the observed ones and then discretize the result into tokens; on the Skytrax dataset with Llama-65B the attack reaches a token accuracy of 88.4 percent, against 22.8 percent for the best baseline. % src: arXiv 2503.09022 abstract
A grey-box variant, in which the adversary knows the base weights and only its own part of the fine-tuned weights, reaches 85.4 percent on Guanaco-65B. % src: arXiv 2503.09022 abstract, Table V
Dong et al.~\cite{dong2025depth} invert both shallow and deep layers of six language models by a two-phase optimization, and report a token F1 of 86.88 on a medical prompt of 4,112 tokens from a middle layer of Llama-3. % src: arXiv 2507.16372 abstract
Thomas et al.~\cite{thomas2025hidden} compute the hidden states of candidate tokens from the vocabulary, match each observed hidden state against them, and recover prompts with nearly perfect accuracy, also when the hidden states are permuted.
Nikolaou et al.~\cite{nikolaou2026language} prove that the map from a prompt to its hidden states is almost surely injective for decoder-only transformers, and their algorithm SipIt recovers the exact prompt one token at a time, demonstrated on GPT-2 Small.
Wan et al.~\cite{wan2024information} study a malicious model provider and train a transformer inverter for the hidden states of deep layers.
With the model in hand, the adversary can therefore compare candidate prompts with the observed hidden states, and recovery is close to exact in some of the evaluated settings.

\paragraph{Attacks with queries.} A weaker adversary holds no weights but can still query the model.
Luo et al.~\cite{luo2025prompt} submit text to the model, collect the hidden states, and train token classifiers on those pairs.
With unlimited queries and auxiliary data, their attacks exceed a token reconstruction accuracy of 90 percent; with a limited query budget and no auxiliary data, the accuracy is typically above 50 percent. % src: arXiv 2503.09291 abstract
The generation attack of Dong et al.~\cite{dong2025depth} likewise trains a sequence-to-sequence inverter on queried pairs of text and hidden states.
In both attacks the model remains callable, and the queries supply the paired examples.
A preprint by Cunningham~\cite{cunningham2026privacy} trains a small classifier directly on about 880 pairs of hidden state and token from Mistral 7B and recovers about 59 percent of the tokens after two layers and about 35 percent after eight, on a test set of 221 pairs. % src: arXiv 2602.16760
That classifier reads each hidden state alone and is evaluated without unpaired matching.

\paragraph{Readout for interpretability.} Interpretability methods also read hidden states, for a different purpose.
The logit lens~\cite{nostalgebraist2020logit} and the tuned lens~\cite{belrose2023eliciting} apply the model's unembedding, directly or after a learned affine map, to a hidden state of an intermediate layer, and Patchscopes~\cite{ghandeharioun2024patchscopes} patches a hidden state into a second forward pass that describes it.
All three methods hold the model.
The two lenses read the model's prediction of the next token, not the token that produced the hidden state, and none of the three is evaluated as an attack.

\section{Inversion of Pooled Text Embeddings}\label{sec:lit:pooled}

A second body of work inverts sentence embeddings, where an encoder pools a whole text into one vector.
Song and Raghunathan~\cite{song2020information} recover between 50 and 70 percent of the words of a sentence from its embedding, as an unordered set, and Pan et al.~\cite{pan2020privacy} infer sensitive attributes, such as identity and health information, from the embeddings of eight language models.
Li et al.~\cite{li2023sentence} move from words to sentences by conditioning a generative decoder on the embedding.
Morris et al.~\cite{morris2023text} refine a hypothesis iteratively: the attack embeds its current guess, compares that embedding with the target, and corrects the guess, and it recovers 92 percent of 32-token inputs exactly from the embeddings of the GTR-base encoder. % src: arXiv 2310.06816 abstract, Table 1
Chen et al.~\cite{chen2024text} extend the attack to multilingual encoders, and Zhang et al.~\cite{zhang2025universal} replace the inverter trained per encoder with one guided search that applies to any encoder; zero-shot there means that no inverter is trained for the encoder, which is still queried.
All of those attacks require access to the encoder, through its weights or through queries; they differ in whether the queries supply paired training examples, score candidate texts during the attack, or both.
The object also differs from a split inference transmission: a pooled embedding is one vector per text, whereas a split model transmits one hidden state per token.

\paragraph{Inversion from outputs.} Prompts can also be reconstructed from what a model returns.
Morris et al.~\cite{morris2024language} train an inverter on the next-token probability vector, and Zhang et al.~\cite{zhang2024extracting} train one on the generated text alone.
Those attacks observe the output of the full model and no hidden state, so they address a different channel.

\section{Cross-Model Representation Alignment}\label{sec:lit:align}

The attacks with weights and the attacks with queries require access that the adversary lacks when the client keeps its blocks private.
Alignment offers another route: the adversary maps the representations of the target model into the space of a model it holds, and inverts there.

\paragraph{Shared structure across models.} That route presumes that independently trained models produce related representations.
Huh et al.~\cite{huh2024platonic} hypothesize that they do: as vision models and language models grow, they measure the distance between inputs in increasingly similar ways.
Moschella et al.~\cite{moschella2023relative} represent each input by its cosine similarities to a fixed set of anchor inputs, and components trained on that relative representation combine across models without further training.
Maiorca et al.~\cite{maiorca2023latent} estimate a direct linear or orthogonal map between two representation spaces in closed form from paired anchors.
For this dissertation the hypothesis is a motivation and not a guarantee: similar structure does not imply that one affine map preserves the identity of each token.

\paragraph{Alignment with little or no supervision.} How many paired examples a map requires has a long history in bilingual word embeddings.
Artetxe et al.~\cite{artetxe2017learning} start from a seed dictionary of 25 word pairs, fit a map, induce a larger dictionary from the mapped embeddings, and repeat; the results are comparable to those of systems with richer bilingual resources. % src: ACL P17-1042 abstract
Conneau et al.~\cite{conneau2018word} and Artetxe et al.~\cite{artetxe2018robust} remove the seed dictionary, the first by adversarial training followed by Procrustes refinement, the second by an initialization from the similarity distributions of the two spaces.
Ormazabal et al.~\cite{ormazabal2019analyzing} examine the limits of that approach: embeddings of two languages learned jointly are more isomorphic than embeddings trained separately and then mapped.
The attack of Chapter~\ref{ch:attack} follows the same pattern, a small paired seed extended by unpaired matching, on the contextual hidden states of language models.
Its matching term adapts the approximate bijective correspondence loss of Murphy et al.~\cite{murphy2022abc}, which learns image representations from set membership, without paired correspondences.

\paragraph{Alignment as an inversion attack.} Three recent attacks apply alignment to sentence embeddings and need no queries.
Huang et al.~\cite{huang2024transferable} train a surrogate encoder to imitate the victim encoder from 8,000 leaked pairs of text and embedding, and train the inverter on the surrogate. % src: ACL 2024.acl-long.230 Sec. 4
Chen et al.~\cite{chen2025algen} fit a linear map from the victim's embedding space to the adversary's from leaked pairs and decode with the adversary's own generator.
A single pair yields a ROUGE-L of about 10; the score rises sharply up to 1,000 pairs and changes little between 2,000 and 8,000. % src: ACL 2025.acl-long.1185 Sec. 5 and ablation
Jha et al.~\cite{jha2025harnessing} remove the pairs entirely: adapters map the embeddings of each encoder into a shared space and are trained on unpaired embeddings with adversarial, reconstruction, cycle-consistency, and vector space preservation losses.
The translated embeddings reach a cosine similarity of up to 0.92 with the true embeddings, and a zero-shot inverter~\cite{zhang2025universal} then extracts part of the text. % src: NeurIPS 2025 paper, results section (0.92); its introduction says 0.96
Those three attacks define the range of supervision, from thousands of pairs to none, and all three act on pooled sentence embeddings.
The matching term of Chapter~\ref{ch:attack} has neither the adversarial objective nor the cycle-consistency losses of Jha et al., so its failure without a leak (Section~\ref{sec:attack:blind}) says nothing about their method.

\section{Defenses}\label{sec:lit:defense}

Defenses for split inference change what the server receives.
Mai et al.~\cite{mai2024split} add noise to the client's token embeddings under local differential privacy and denoise the model's output on the client.
A second family hides the hidden states behind a secret permutation, as in PermLLM, STIP, and Centaur~\cite{zheng2024permllm,yuan2023secure,luo2025centaur}.
Thomas et al.~\cite{thomas2025hidden} show that all three are insecure, since the matching of hidden states to vocabulary candidates can be made invariant to the permutation, and propose Cascade, which divides the token sequence among servers that must not collude.
Dong et al.~\cite{dong2025depth} test four defenses, among them quantization and differential privacy, and report that none fully prevents their inversion.
The evaluations of Thomas et al. and Dong et al. assume an adversary that holds the model's weights.
This dissertation evaluates undefended hidden states, so its results establish no break of any of those defenses.

\section{Position of the Dissertation}\label{sec:lit:position}

Table~\ref{tab:lit:position} summarizes the access that each group of attacks assumes.
The attacks on language models listed in the table fall into two groups.
In the first, the adversary holds or queries the model that produced the per-token hidden states.
In the second, the attacks act on pooled sentence embeddings, with queries to the encoder or, through alignment, without them.

\begin{table}[ht]
\caption{The listed attacks hold or query the target model, or act on pooled embeddings; the dissertation fits a map from per-token hidden states without the target model's weights and without queries.
The settings of the dissertation's attack were selected on additional prompts with known text.}
\label{tab:lit:position}
\centering
\ssp\small
\begin{tabular}{@{}>{\raggedright\arraybackslash}p{0.30\textwidth}lcc>{\raggedright\arraybackslash}p{0.22\textwidth}@{}}
\toprule
Attack & Representation & Weights & Queries & Paired examples \\
\midrule
Qu, Thomas, Nikolaou et al.~\cite{qu2025prompt,thomas2025hidden,nikolaou2026language} & per token & yes & no & none \\
Luo, Dong et al.~\cite{luo2025prompt,dong2025depth} & per token & no & yes & from queries \\
Morris et al.~\cite{morris2023text} & pooled & no & yes & from queries \\
Huang et al.~\cite{huang2024transferable} & pooled & no & no & 8,000 leaked \\
Chen et al.~\cite{chen2025algen} & pooled & no & no & 1 to 8,000 leaked \\
Jha et al.~\cite{jha2025harnessing} & pooled & no & no & none \\
This dissertation & per token & no & no & \vf{90} leaked rows \\
\bottomrule
\end{tabular}
\end{table}

To the author's knowledge, prior work has not evaluated query-free prompt reconstruction from per-token hidden states of a language model by fitting a map on a counted leak together with unpaired matching.
Each part of that setting has precedent: per-token hidden states are inverted with weights or queries, sentence embeddings are inverted from a few leaked pairs, and embedding spaces are aligned without pairs.
The dissertation combines them, and it measures what the unpaired matching adds to a fixed leak.
Two qualifications bound the statement.
A row is a token group of a leaked text span, not an isolated pair of hidden state and token, and the settings of the attack were selected on prompts with known text outside the leak.
The adversary of Chapter~\ref{ch:attack} also holds the hub trained on public models, the tokenizer of the target model, and unpaired data from the domain under attack.
